\documentclass[11pt]{article}

\usepackage[T1]{fontenc}
\usepackage[utf8]{inputenc}
\usepackage[letterpaper,margin=1in]{geometry}
\usepackage{amsmath}
\usepackage{booktabs}
\usepackage{graphicx}
\usepackage{xcolor}
\usepackage{listings}
\usepackage{hyperref}

\definecolor{linknavy}{rgb}{0.15,0.25,0.55}
\hypersetup{
  colorlinks=true,
  linkcolor=linknavy,
  citecolor=linknavy,
  urlcolor=linknavy,
  breaklinks=true,
  pdftitle={fleetsim: a validated discrete-event simulator for machine-learning
            accelerator fleets},
  pdfauthor={Felipe Valencia},
}

% Listings are configured but the paper currently sets no code blocks; the
% style is here so that adding one does not change the document's look.
\lstset{
  basicstyle=\ttfamily\small,
  breaklines=true,
  columns=fullflexible,
  frame=single,
  framerule=0.3pt,
  rulecolor=\color[gray]{0.7},
  showstringspaces=false,
}

% Tables in the fragments are dense: more vertical air, less horizontal
% padding, so the widest ones fit the text block.
\renewcommand{\arraystretch}{1.12}
\setlength{\tabcolsep}{5pt}

% The paper is mostly tables. Loosen LaTeX's default float limits so a large
% table lands near the text that discusses it instead of being deferred to
% the end of the document.
\setcounter{topnumber}{3}
\setcounter{bottomnumber}{2}
\setcounter{totalnumber}{5}
\renewcommand{\topfraction}{0.92}
\renewcommand{\bottomfraction}{0.6}
\renewcommand{\textfraction}{0.06}
\renewcommand{\floatpagefraction}{0.7}

% Several paragraphs contain unbreakable tokens (long URLs, full-precision
% floats, shell commands). Let TeX stretch a line as a last resort rather
% than run it into the margin, and let URLs break at any character.
\setlength{\emergencystretch}{2em}
\makeatletter
\g@addto@macro{\UrlBreaks}{\UrlOrds}
\makeatother

\title{\texttt{fleetsim}: a validated discrete-event simulator for
       machine-learning accelerator fleets}

\author{Felipe Valencia\\[2pt]
  \normalsize\url{https://github.com/felivalencia3/mini-compute-simulator}}

\date{\today}

\begin{document}

\maketitle

\begin{abstract}
\noindent
\texttt{fleetsim} is an open-source discrete-event simulator for fleets of
machine-learning accelerators. It models a hierarchy of clusters, pods, racks
and nodes; gang scheduling, where a job needs all its chips at once or none;
preemption priced in the training work it destroys; node failures, repair and
maintenance drains; and a traffic generator whose parameters come from published
cluster characterizations. Existing open tools sit at other \emph{altitudes} ---
altitude meaning what one simulated run represents: a single cluster with a flat
GPU count, or one job's network traffic with no queue in front of it --- and the
production fleet managers at this altitude ship nothing anyone outside can run.

The central claim is validation against published measurements of real
clusters. On the four Helios GPU clusters, \texttt{fleetsim} reproduces the
published ratio between first-in-first-out (FIFO) and shortest-job-first mean
completion times inside the stated tolerance bands on all four, at a
four-cluster mean absolute error of 3.4\,\%. The absolute mean FIFO completion
time on Saturn, the busiest of the four, lands within 0.01\,\% of published ---
a coincidence, not a precision claim, because reordering jobs inside the trace's
tied submit seconds moves that same quantity by 17\,\%. What the absolute
comparison supports is agreement within $\pm 14$\,\% across all four clusters,
never a single cell. Two further qualifiers travel with the result: the bands
are wide for stated reasons --- a deliberately bad placement policy passes them
--- and every number holds only under three named modelling choices. We also
list the published quantities \texttt{fleetsim} deliberately does \emph{not}
reproduce.

\medskip
\noindent\textbf{Availability.} \texttt{fleetsim} is Apache-2.0 licensed at\\
\url{https://github.com/felivalencia3/mini-compute-simulator}\\
The offline checks are one command, \texttt{pytest tests validation -q}. The
Helios replay is environment-gated: it downloads the 36\,MB public trace and
replays four clusters under two schedulers in roughly four minutes
(250--266\,s across the repository's own timings, on a laptop).
\end{abstract}

\section{Introduction}
\label{sec:intro}

A fleet of machine-learning accelerators is a large capital asset that behaves
like a queueing system. The chips are bought in bulk, wired into a fixed
hierarchy, and handed out --- minute by minute, for months --- by one program:
the scheduler, which decides which queued job runs now and which waits.

That decision has more leverage than it looks. On a 2{,}048-chip fleet --- the
shipped \texttt{examples/01\_minimal} scenario, not one of the case studies in
\S\ref{sec:cases} --- replacing strict first-in-first-out (FIFO) queueing with a
preempting tiered-priority policy cuts the mean number of waiting evaluation
jobs from about 11 to about 1.3, and their 99th-percentile queue wait from about
5\,h to about 31\,min. Same hardware, same arrivals, same seed; a different
ordering rule. (Those are the figures the example's own README records, to the
precision it records them.)

You cannot run that experiment on the production fleet: an A/B test costs weeks
of accelerator time, and traces record what one scheduler did, never what another
would have done.

\texttt{fleetsim} is an open-source discrete-event simulator that answers such
questions at \emph{fleet altitude}. Altitude here means what one simulated run
represents. One \texttt{fleetsim} run is weeks of a whole fleet: thousands of
jobs arriving, queueing, running, being preempted and occasionally dying with
their nodes. A network simulator's run, by contrast, is one already-placed job
executing a single training step. Given a fleet in YAML, arriving jobs and a
failure model, \texttt{fleetsim} reports the occupancy, queue waits, preemption
rates and goodput a policy delivers --- four quantities
\S\ref{sec:threenumbers} and \S\ref{sec:failures} define, and whose differences
turn out to matter more than any one of them alone. It models no packets,
collectives or kernels --- per-job speed is one analytical multiplier ---
because its questions concern thousands of jobs contending for one machine, not
one job's step time.

We claim three contributions, in ascending order of importance.

\begin{enumerate}
\item \textbf{A fleet-altitude model} (\S\ref{sec:models}): one capacity tree
from metro down to node, spanning NVIDIA-style clusters and TPU-style pods, with
gang scheduling, locality constraints, preemption cost, quota, failures and
drains in one engine; the simulators we surveyed each model a subset
(\S\ref{sec:related}).
\item \textbf{A generative traffic model} (\S\ref{sec:traffic}): arrivals, job
sizes and durations from published cluster characterizations rather than a single
Poisson assumption, each parameter recorded with its source --- including those
that are order-of-magnitude priors rather than fitted constants.
\item \textbf{A validation suite} (\S\ref{sec:validation}) that replays a real
production trace --- the four Helios clusters~\cite{helios2021} --- and
reproduces the published FIFO-versus-shortest-job-first completion-time ratios to
a mean absolute error of 3.4\,\%, under three stated modelling choices we restate
wherever the result appears. It adds closed-form queueing checks, a list of
published numbers \texttt{fleetsim} deliberately does \emph{not} reproduce, and
one mechanism we had misdiagnosed (\S\ref{sec:replay-found}).
\end{enumerate}

\S\ref{sec:background} primes readers new to this infrastructure;
\S\ref{sec:models}--\ref{sec:traffic} give the model and traffic generator,
\S\ref{sec:validation} the validation, \S\ref{sec:cases}--\ref{sec:limitations}
the case studies, implementation, related work and limitations.

\section{Background: what a fleet is, and why scheduling it is hard}
\label{sec:background}

\subsection{The physical machine}

The unit of capacity is a \emph{chip}: one GPU or one TPU. Eight typically share
a server chassis and a fast internal switch --- a \emph{node}. Tens of nodes make
a rack or scalable unit, and tens of those a \emph{pod}: 4{,}096 GPUs in the H100
configurations we model, and 4{,}096 chips in a TPU v4 pod~\cite{tpuv4isca2023}.
Pods join into \emph{clusters}, then datacenters and \emph{metros}.

Every level is a bandwidth cliff, which is why locality matters. Llama~3's
cluster was eight pods of 3{,}072 GPUs at 1:7 oversubscription, and the 405B run
took 16{,}384 of those GPUs~\cite{llama3}. The 1:7 figure means seven times less
bandwidth between pods than within one, and a job straddling such a boundary
runs slower for its whole life.

\subsection{Gang scheduling and fragmentation}

A distributed training job is not a web service: its chips run one computation in
lockstep, synchronizing every step, so it needs \emph{all} of them \emph{at once}
for its whole life. Half a training job does not do half
the work; it does none. Allocation is therefore all-or-nothing --- \emph{gang
scheduling}. The analogy is a restaurant booking for twelve: a table for eight
plus a table for four across the street does not seat twelve.

The gang is also the unit of failure: if any node holding part of it dies, the
whole job dies. And gangs create \emph{fragmentation}, free capacity that cannot
be used: a fleet can have 200 chips free and still not start a 64-chip job,
because those chips sit in ones and twos across many nodes. Philly
measurements attribute 59--98\,\% of multi-GPU queueing delay to fragmentation
rather than fair-share policy~\cite{philly2019}; reproducing that attribution is
itself an explicit non-goal (\S\ref{sec:anti-goals}).

Three kinds of job share the machine, and they are not variations on a theme.
A \emph{pretraining} run trains a model from scratch: thousands of chips held in
lockstep for weeks. A \emph{finetune} adapts an already-trained model to a
narrower task: tens of chips for hours. An \emph{evaluation} scores an existing
model on a test set: one or two chips for minutes. All three arrive into the
same queue, and they differ by about four orders of magnitude in both size and
duration.

That disparity shows up in every trace. Over 82\,\% of Philly's jobs use one
GPU~\cite{philly2019}, and 70.8\,\% of jobs on the busiest Helios cluster we
replay are single-GPU. Acme measured evaluations at about 93\,\% of job
\emph{count} and under 1\,\% of GPU-\emph{time}, with pretraining at 1--3\,\% of
count and 70--94\,\% of GPU-time~\cite{acme2024}. The tiny jobs strand exactly
the capacity the enormous ones need.

This is the \emph{mice-and-hogs} shape: a swarm of one- and two-chip jobs (the
mice) that dominate the job count, and a handful of thousand-chip jobs (the
hogs) that dominate the chip-hours. Almost every metric in this paper has to be
reported twice because of it. Count one job as one job and you are describing
the mice; weight by chip-hours and you are describing the hogs.

A fleet is also shared. Each \emph{tenant} --- a team or project with its own
budget --- is granted a \emph{quota}, a promised share of the chips the
scheduler is expected to honour before it hands anything to anyone else. Quota
is a promise about capacity, not a physical partition: the chips themselves stay
in one pool.

\subsection{Three numbers that are not the same number}
\label{sec:threenumbers}

One term first, because the third definition below depends on it. A long
training job periodically writes its entire state to storage --- a
\emph{checkpoint} --- so that a job killed by a failure or a preemption resumes
from the last saved state rather than from the beginning. Saving costs the job
real time (about a minute, on \texttt{fleetsim}'s default one-hour interval),
and everything computed since the last checkpoint is lost when the job dies. The
analogy is a document editor that autosaves every hour: a crash costs you at
most an hour, and the autosaves themselves cost you a little throughput.

Conflating utilization metrics is the classic error here, so \texttt{fleetsim}
defines three and always reports all three. \emph{Allocation rate} is allocated
chip-time over \emph{total} chip-time, including chips that are broken or
drained. \emph{Occupancy} is allocated chip-time over \emph{healthy} chip-time.
\emph{Goodput} is the fraction of allocated chip-time that produced surviving
work, excluding work lost back to the last checkpoint, restart latency and
checkpoint-save overhead.

A worked example. Of 1{,}000 chips, fifty are down for repair, so 950 are
healthy, and all 950 are allocated. Allocation rate is $950/1000 = 0.95$;
occupancy is $950/950 = 1.00$, and the gap is exactly the capacity lost to
failures --- which is why occupancy alone makes a fleet with broken hardware look
perfect. Now suppose a third of the allocated chip-time goes on redoing work lost
at the last interruption and on restarts: goodput is $0.67$, and the share of the
machine doing useful work is $0.95 \times 0.67 \approx 0.64$. Fully occupied, a
third of its work wasted. Philly clusters were in fact measured at roughly
100\,\% allocation with about 52\,\% hardware utilization~\cite{philly2019}.

\subsection{Failures and preemption}
\label{sec:failures}

At small scale hardware failure is a nuisance; at fleet scale it is a design
constraint, because a gang fails when \emph{any} of its nodes fails. A gang of
$N$ nodes therefore breaks $N$ times as often as one node does.

Two published anchors bracket the numbers. Meta measured node mean time between
failures (MTBF) at about 42 node-days~\cite{metareliability2024}: one node in
the fleet breaks about every 42 days. Llama~3's 405B run, 16{,}384 GPUs and so
about 2{,}048 nodes, saw roughly one unplanned interruption every 3.1
hours~\cite{llama3}. Scaling that measurement, a gang eight times larger ---
16{,}384 nodes, the size our frontier case study uses --- would be interrupted
eight times as often, about every 23 minutes.

The two anchors do not agree with each other, and we say so rather than picking
the convenient one. Applying the 42-node-day MTBF directly to Llama~3's 2{,}048
nodes predicts an interruption every 30 minutes, roughly six times more often
than Llama~3 actually reported. The MTBF checks out against Meta's own smaller
figure (42 node-days over the 128 nodes of a 1{,}024-GPU job gives 7.9 hours,
which is what Meta measured) and not against Llama~3's. \texttt{fleetsim}
therefore ships 42 node-days as a configurable default, not as a fitted
constant, and the 23-minute figure quoted later comes from scaling Llama~3's
measurement, not from the MTBF.

Preemption --- stopping a running job to give its chips to a more important one
--- costs the same way: a preempted training job does not pause, it loses
everything since its last checkpoint and pays about 15 minutes of restart
overhead~\cite{metareliability2024}. Under a one-hour checkpoint interval that is
up to an hour of work per chip in the gang. This is how aggressive
preemption raises occupancy while destroying goodput.

\section{What fleetsim models}
\label{sec:models}

\subsection{One tree for two ecosystems}

Capacity is a tree of \emph{domains} from metro down to node, all one
\texttt{Domain} class. Leaves are nodes carrying a chip count, a chip type, a
health state (\texttt{HEALTHY}, \texttt{DRAINING}, \texttt{FAILED},
\texttt{MAINTENANCE}) and a failure-rate multiplier. Interior domains carry chip
counters maintained incrementally, making ``how much free capacity is under this
domain'' an $O(1)$ query.

One tree covers both ecosystems because \emph{level names are configuration, not
code}: a cluster declares \texttt{[cluster, pod, su, node]} for an H100 SuperPOD
or \texttt{[cluster, pod, cube, host]} for a TPU v5p pod. The scheduler sees
domains, capacities and constraints, never vendors.

\subsection{The allocation rule, and why it is the hardest call}

Jobs are allocated chip-granularly \emph{within} a single node and in whole-node
multiples \emph{above} one node: a 2-chip evaluation may share a node, a 64-chip
job asks for eight entire nodes. The reason is workload shape ---
\texttt{fleetsim}'s shipped class mix puts about 90\,\% of arrivals in 1--8-chip
evaluations, following the Acme characterization~\cite{acme2024}, which measured
the corresponding share at about 93\,\%. Were the node atomic,
a 1-chip evaluation would strand seven GPUs and corrupt both occupancy and the
mice-and-hogs story; rounding those jobs up would falsify the workload.

The rule has one load-bearing consequence, the mechanism behind our main
validation finding: \emph{a leaf with any owner at all is invisible to every
whole-node request}. A node with one chip in use has seven free chips no
multi-node gang can claim, so free capacity accumulates as 1--7-chip remainders.
That whole-node rule is an approximation we measured: in the Helios September
window the
trace's own node count satisfies $\texttt{node\_num} = \lceil
\texttt{gpu\_num}/8 \rceil$ for 99.66\,\% of multi-GPU jobs, and all 225
exceptions used \emph{more} nodes than the ceiling, never fewer --- the half the
model depends on.

An allocation is atomic, immutable once granted, and the unit of failure. A gang
may name a level it must fit \texttt{within}, and may be \emph{segmented} into
whole-node blocks that each sit inside one domain at that level.

\subsection{The engine}

Time is an \texttt{int64} count of microseconds, which makes tie-breaking exact.
The engine is event-driven rather than tick-driven because
state changes are sparse and unevenly spaced: a fleet can go minutes without one,
then process hundreds of preemptions at a single instant.

Events sit in a heap keyed by \texttt{(time, type, seq)}. The event type is not
decoration --- its integer value \emph{is} the same-timestamp ordering rank,
arranged so every repair, completion, failure and preemption at time $t$ is
processed before the scheduler wakes at $t$, so it always sees settled state.
Scheduler wakes are then
\emph{coalesced}: events set a dirty flag and ensure at most one pending wake at
the next round boundary, so the scheduler runs once per round (60\,s by default)
rather than once per event.

\subsection{Schedulers and placement}

Ordering and placement are separate axes~\cite{blox2024}: a scheduler sees an
immutable view of the cluster and returns \texttt{Place} and \texttt{Preempt}
intents, which the engine validates and applies~\cite{batsim2016}.

Four schedulers ship. \texttt{fifo} serves the queue in arrival order, with
strict head-of-line blocking by default: the head of the queue holds up
everything behind it until it can run. \texttt{tiered\_priority} implements
Borg-style priority bands~\cite{borg2015,borgtng2020} with requeue preemption.
\texttt{sjf} runs the shortest job first. \texttt{easy\_backfill} adds
\emph{backfilling}~\cite{easybackfill1995}: letting a small job at the back of
the queue jump ahead into a hole the head job cannot use, on the condition that
doing so provably does not delay the head. That condition is computed from each
job's declared runtime estimate, and users routinely over-estimate, so a
backfill scheduler must behave sanely when the estimate is wrong. Plugins
register through a Python entry point.

Orthogonally, four placement policies choose \emph{which} fitting nodes to take.
\texttt{first\_fit} (the default) takes the first nodes that fit, scanning by
ascending node identifier. \texttt{best\_fit} takes the nodes with the least
free space that still fits, so partly-used nodes fill up before empty ones are
opened --- \emph{tightest-fit} packing, a term used throughout the rest of the
paper. \texttt{consolidate} does the same at the sub-node level and additionally
minimises how many parent domains a multi-node gang touches; on a single-level
fleet it is bit-identical to \texttt{best\_fit}. \texttt{spread} takes the
emptiest nodes (worst fit) --- a deliberate anti-policy, our control arm.

\subsection{Failures, repair and maintenance}

Failures are sampled by \emph{aggregate} rate --- one exponential draw at the
fleet-wide rate, then a victim node --- not by per-node timers, so a job's mean
time to interruption scales inversely with its size for free.

Defaults are sourced, not invented: node MTBF 42 node-days, automatic repair
uniform on 60--180 minutes, 10\,\% of failures needing manual repair over 1--3
days, one maintenance drain per node-month, a one-hour drain grace window, and a
failure-cause mix of $\sim$60\,\% GPU/HBM, $\sim$10\,\% network and
$\sim$13\,\% software, with the remaining 17\,\% a residual ``other'' bucket
\texttt{fleetsim} adds so the mix sums to one. The MTBF comes from Meta's
measurements~\cite{metareliability2024}, the three approximate cause shares from
Llama~3~\cite{llama3}, and the drain rate from Borg~\cite{borg2015}. The repair
distributions are the weakest-sourced of the set: the repository attributes them
to AIReSim~\cite{aireesim}, a reliability simulator it names but does not pin to
a citable artefact, so a reader cannot check them the way the others can. All
are configuration values, none universal. Maintenance is a separate stream: a
\texttt{DRAINING} node takes no new placements, and its residents are requeued
when the grace window expires.

\subsection{Metrics}

Table~\ref{tab:metrics} states the definitions. State metrics are time-weighted
integrals updated on every state change --- exact, not sampled --- and
accumulated as integer chip-microseconds; discrete outcomes are event-sourced.

\begin{table}[t]
\centering
\caption{Metric definitions. ``Allocated'' means chips held by a running job;
``healthy'' excludes nodes that are failed, draining or under maintenance.
Integrals are exact integer chip-microseconds internally, reported as
chip-seconds.}
\label{tab:metrics}
\small
\begin{tabular}{@{}lp{0.36\linewidth}p{0.30\linewidth}@{}}
\toprule
Metric & Definition & What it tells you \\
\midrule
Allocation rate & $\int \text{allocated}\,dt \,/\, \int \text{total}\,dt$ &
Share of the purchased fleet handed out \\
Occupancy & $\int \text{allocated}\,dt \,/\, \int \text{healthy}\,dt$ &
Occupancy minus allocation rate is capacity lost to failures and drains \\
Goodput & $\sum \text{productive chip-s} \,/\, \int \text{allocated}\,dt$ &
Excludes work lost since the last checkpoint, restart latency and
checkpoint-save overhead \\
ETTR (per job) & Effective training time ratio: productive seconds / wall
seconds holding an allocation &
Goodput seen from one job's side; counts restart overhead and preemption grace
in the denominator \\
Fragmentation index & $1 - \text{largest placeable} / \text{total free}$ at a
level & $0$ when all free capacity is usable by one job \\
Stranded whole nodes & healthy leaves with $0 < \text{free} < \text{chips}$ &
Free capacity no whole-node gang can claim \\
\bottomrule
\end{tabular}
\end{table}

Two conventions matter. Metrics are given over the full run and over a
steady-state window (the middle 80\,\% of the horizon by default), and every
distributional metric is reported job-weighted \emph{and} chip-hour-weighted ---
a job-weighted mean over a mice-dominated workload describes the mice.
Best-effort jobs are excluded from wait distributions. \emph{Best-effort} is the
lowest, always-preemptible priority band, and \texttt{fleetsim} submits it from
a \emph{closed loop}: a source that keeps a fixed number of jobs in flight and
launches a replacement the instant one finishes. A closed loop by construction
always has work queued, so its queue length says nothing about the fleet and its
mean waiting time is undefined. Best-effort work is therefore reported by its
goodput instead.

\section{Modelling the traffic}
\label{sec:traffic}

Much of what a fleet simulator tells you depends on what arrives at it. The
textbook default is Poisson arrivals: jobs appearing independently of one another
at a constant average rate, with exponentially distributed gaps. Poisson traffic
is smooth in a specific sense --- average it over a longer window and its
relative variability shrinks predictably. Measured computer traffic does not do
that. It is bursty at every timescale you look at, and averaging does not smooth
it out \cite{paxsonfloyd1995}. A study built on Poisson arrivals understates
queue variance, and understates it most exactly where the answer matters, near
saturation. fleetsim therefore runs one arrival process per class and superposes
them (Table~\ref{tab:arrivals}).

\begin{table}[t]
\centering\small
\caption{Arrival process per job class, with shipped defaults. Pretraining is
the one surviving Poisson case: rare planned launches really are independent,
human-initiated events.}
\label{tab:arrivals}
\begin{tabular}{lll}
\toprule
Class & Process & Defaults \\
\midrule
Pretrain    & homogeneous Poisson              & \texttt{rate\_per\_week: 2} \\
Finetune    & MMPP-2 $\times$ diurnal envelope & \texttt{rate\_ratio: 4}, \texttt{burst\_frac: 0.25}, \texttt{switch\_tau: 2d} \\
Eval        & Hawkes on an NHPP baseline       & \texttt{branching: 0.4}, \texttt{kernel\_tau: 15m} \\
Inference   & service replica curve            & --- \\
Best-effort & closed-loop backlog              & \texttt{concurrency: 4}, \texttt{think\_time: fixed[0s]} \\
Retry storm & near-critical Hawkes (preset)    & \texttt{branching: 0.9}, \texttt{kernel\_tau: 2m} \\
\bottomrule
\end{tabular}
\end{table}

The shared day-and-week machinery is a non-homogeneous Poisson process whose log
rate is a sum of harmonics at 24\,h and 168\,h. The exponential keeps the rate
positive, and the constant term is derived rather than configured, so changing
the shape of a day does not change the amount of work. The shipped daily preset
is a two-harmonic least-squares fit to \texttt{fleetsim}'s own v0.1 step curve,
which was itself built from the Helios description of a night trough with dips
at 12:00 and 18:00 \cite{helios2021}. The provenance matters because the fit
loses something: two harmonics smooth those two dips away, leaving a trough near
$0.59\times$ at 03:00 and a single peak near $1.27\times$ at 10:15. The shipped
preset is therefore continuous with the previous release's behaviour rather than
a fit to a published curve, and it no longer reproduces the midday and evening
dips Helios reports. It is sampled by Lewis--Shedler thinning --- propose
arrivals at a constant upper-bound rate, then keep each with probability equal
to the true rate over that bound.

Finetuning sits on a two-state Markov-modulated Poisson process: a hidden switch
between a quiet and a busy regime, reproducing crunch weeks that a single rate
curve cannot. It is configured in observable terms --- mean rate, burst fraction,
switch correlation time, rate ratio --- and solved to state rates in closed form.
Those rates must be scaled to the \emph{peak} of the diurnal envelope before
thinning; left unscaled they silently generate about 60\,\% of the configured
load.

Evaluation traffic uses a Hawkes process. Self-exciting means one submission
makes further submissions more likely for a while afterwards: a failed eval is
re-run, a user iterates. Configuration exposes the branching ratio --- mean
children per arrival, stationary below 1 --- and the kernel decay time; sampling
uses Ogata thinning, sequential but $O(1)$ per event because the excitation term
decays exponentially.

Sizes are per-class distributions over powers of two. Evals are 55\,\%
single-chip: inside the 50--90\,\% single-GPU range Helios reports, and
deliberately below Philly's finding that over 82\,\% of jobs are single-GPU
\cite{philly2019}. TPU clusters use the published histogram of \emph{slice}
sizes \cite{tpuv4isca2023} --- a TPU slice being a contiguous geometric block of
chips carved out of a pod, a different thing from the quota slices of
\S\ref{sec:helios}. Uniform-over-exponents sizing remains available but is not
the default, since it overweights giants by 5--10$\times$ and fabricates
queueing.

Job durations are drawn from a lognormal shape, configured by naming its
quantiles (median and 90th percentile) rather than its abstract parameters.
Pretraining additionally gets a \emph{heavy tail}: a small probability of a run
lasting far longer than the median, decaying as a power law rather than dying
off exponentially. How heavy a tail is, is summarised by its \emph{tail index};
a smaller index means a heavier tail. Pretraining gets the only heavy tail
because resource-hours are size $\times$ duration, and a product of two
quantities is as heavy-tailed as the heavier of them --- it inherits the
\emph{smaller} of the two indices. Attaching the tail to the one class with
genuinely open-ended runs is therefore intended to suffice for the Borg-like
resource-hour tail \cite{borgtng2020}, without fabricating implausibly long
evaluation jobs. Intended, not confirmed: the distributional self-check that
would measure the generated chip-hour tail is a documented target in the traffic
specification, not a shipped test.

Tenants come from a \emph{finite} Zipf over ranks. An exponent of 1.19
reproduces Alibaba PAI's top 5\,\% of users submitting about 77\,\% of
instances \cite{alibabapai2022}, at a population of 1{,}300 tenants; the shipped
default is 1.2. The class mix follows Acme \cite{acme2024}.

These parameters are sourced but not all fitted, and the repository says so: the
Markov-modulated and Hawkes defaults are order-of-magnitude priors borrowed from
grid and LLM-serving analogies, not constants fitted to an ML fleet. There is no
explicit long-range-dependence generator, tenants are marks on jobs rather than
agents, and within a class size and duration are drawn independently.

\section{Validation}
\label{sec:validation}

\subsection{What ``validated'' means here, and the two kinds of check}
\label{sec:val-kinds}

A simulator can be internally consistent and still be wrong about the world.
Validation here means comparing fleetsim's output against numbers a published
study measured on a real cluster, and two kinds of comparison are possible. A
\textbf{policy-effect check} runs two scheduling policies over the same trace and
compares the \emph{ratio} between them against the ratio the paper reported. A
\textbf{distribution-match check} asks whether a converted trace reproduces a
published shape, such as the fraction of jobs that ended successfully; that
mostly tests the converter rather than the scheduler, and in the Philly case
involves no simulation at all.

The ratio is the more robust test, and the reason deserves two extra sentences.
Suppose the fleet we reconstructed is 10\,\% smaller than the real one. Both
policies then run on a fleet 10\,\% too small, both produce waits that are too
long, and much of that shared error divides out when one is taken over the other.
An absolute number has nothing to divide it out; it carries every reconstruction
error straight through.

Replays never match exactly, for structural reasons. Take our own replay first.
Helios publishes its analysis window --- 2020-09-01 to 2020-09-26 --- and we
replay exactly that, so the window is not a source of error here. Two others
remain. Timestamps carry no timezone, so a daily cycle cannot be aligned with
confidence. And the original scheduler made placement decisions nobody recorded,
which is the residual \S\ref{sec:replay-found} spends most of its length on.

Where a released trace and its paper's analysis window differ, the problem is
worse. Philly released 117{,}325 jobs over 137 days while the paper analysed
96{,}260 over roughly 75, and that window is not precisely published, so no
reconstruction of it can be checked. The tolerance bands below are sized to
absorb all of those residuals; they are not calibrated uncertainty intervals.

\subsection{Replaying a real trace: the Helios experiment}
\label{sec:helios}

The flagship replay targets Hu et al.'s Helios study \cite{helios2021}, which
characterizes four production GPU clusters --- Venus, Earth, Saturn and Uranus
--- over 2020-09-01 to 2020-09-26, GPU jobs only. The trace is a
36{,}437{,}672-byte archive (343{,}069{,}364 bytes uncompressed) under CC-BY-4.0;
fleetsim fetches it, verifies it by size, and replays all four clusters in about
four minutes.

Each cluster is partitioned into \emph{virtual clusters}: quota slices, each a
group's private share of the machines, scheduled independently. The reference
simulator runs one scheduler per slice, whereas fleetsim schedules a single
global pool, so the replay is a harness rather than an engine change --- one
simulation per slice, on a fleet sized to that slice's node count at 8 GPUs per
node, aggregated job-weighted to cluster totals. The converter writes each job's
recorded duration into both its duration and its runtime estimate, making
shortest-job-first an \emph{oracle}: a perfect estimate, and so a strict upper
bound on what a duration-\emph{predicting} policy could achieve.

Every claim below holds only under three modelling choices, all options a reader
could disagree with: a strict blocking scan, September-max per-slice sizing, and
\texttt{consolidate} placement. Two of them were not chosen but forced ---
Table~\ref{tab:corrections} records the assumptions the real data disproved. A
strict scan means the queue head blocks everyone behind it until it can run,
which is what head-of-line blocking means; both policies share the scan mode, so
shortest-job-first's advantage remains purely its ordering. September-max sizing
drops zero jobs, and the test asserts that, but it is not unambiguously more
faithful and the repository says so: peak quota available from day one
over-provisions early in the month and under-counts early-month queueing. The
converter's own default remains the snapshot.

\begin{table}[t]
\centering\small
\caption{Two modelling assumptions the validation plan made before the data was
touched, and what the real trace showed.}
\label{tab:corrections}
\begin{tabular}{p{0.14\textwidth}p{0.30\textwidth}p{0.46\textwidth}}
\toprule
Choice & The plan assumed & What the trace showed \\
\midrule
Scan mode &
Non-blocking: a job the scheduler cannot place is skipped and later jobs
proceed. &
The reference blocks. Non-blocking collapses FIFO's queueing --- ratios near
$1.4\times$ with the wrong cross-cluster ordering. A strict scan flips the
result from non-reproducing to reproducing. \\
Per-slice capacity &
A fixed quota snapshot taken on 1 September. &
Quota drifts mid-month: one Uranus slice grows 208 $\to$ 328 GPUs, another
spins up 0 $\to$ 416, and 3{,}117 jobs sit in slices with no 1 September quota
at all, which a snapshot silently drops. It also over-congests Uranus (ratio
$3.41\times$), breaking the rank and the queuing-share ordering. \\
\bottomrule
\end{tabular}
\end{table}

Table~\ref{tab:helios-ratios} gives the result: all eight ratios in band. The
quantity being compared is the job completion time, or JCT --- the wall-clock
time from when a job was submitted to when it finished, so queue wait plus run
time. The test also asserts that no job was dropped, that every job reached a
terminal state, that Saturn has the highest and Uranus the lowest JCT ratio, and
that the queuing share of JCT orders Saturn $>$ Venus $>$ Earth $>$ Uranus in
fleetsim (90.1, 79.1, 73.0, 42.2\,\%) as in the published table (89.7, 81.8,
69.3, 42.5\,\%). It runs in roughly four minutes on a laptop against the cached
trace: the repository's own timings are 250\,s, ``about 260\,s'' and 266\,s in
three different files, which is the accuracy a wall-clock number of this kind
deserves. The bands are wide on purpose, and were deliberately \emph{not}
tightened even though every JCT ratio now lands within 5\,\% of published, for
the two reasons in \S\ref{sec:val-kinds}. Nor is the queuing column an exact
reproduction: Earth measures $10.67\times$ against a published $16.4\times$,
about 35\,\% low, and passes only because the band runs to $25\times$. The band
does have teeth --- SJF replayed against itself gives 1.0 and fails it.

\begin{table}[t]
\centering\small
\caption{Helios replay, four clusters, real trace, shipped configuration (strict
scan, September-max sizing, \texttt{consolidate} placement). The
\texttt{fleetsim} column is the claim. The \texttt{first\_fit} column is the same
replay under the engine's \emph{default} placer, shown so the size of the
placement effect is visible rather than asserted. The \emph{In band} columns
judge the \texttt{fleetsim} column only. Bands are $[1.3, 8]\times$ on the
completion-time ratio and $[3, 25]\times$ on the queuing ratio; $\dagger$ marks
the one \texttt{first\_fit} value that falls outside its band, the deviation
\S\ref{sec:replay-found} diagnoses. The mean absolute errors quoted in the text
(3.4\,\% for \texttt{fleetsim}, 27.4\,\% for \texttt{first\_fit}) are computed on
full-precision ratios; the rounded values shown here give 3.2\,\% and 27.5\,\%.}
\label{tab:helios-ratios}
\begin{tabular}{lrrcrrrcr}
\toprule
& \multicolumn{4}{c}{FIFO/SJF mean job-completion-time ratio} & \multicolumn{4}{c}{FIFO/SJF mean-queuing ratio} \\
\cmidrule(lr){2-5}\cmidrule(lr){6-9}
Cluster & Published & fleetsim & In band & \texttt{first\_fit} & Published & fleetsim & In band & \texttt{first\_fit} \\
\midrule
Saturn & 6.59 & 6.87 & yes & 8.75$^\dagger$ & 18.5 & 19.55 & yes & 22.91 \\
Venus  & 3.07 & 3.21 & yes & 4.21           & 5.68 & 7.78  & yes & 9.97  \\
Earth  & 2.87 & 2.95 & yes & 2.11           & 16.4 & 10.67 & yes & 5.73  \\
Uranus & 1.49 & 1.51 & yes & 1.69           & 4.51 & 5.08  & yes & 6.10  \\
\bottomrule
\end{tabular}
\end{table}

\subsection{What the replay found}
\label{sec:replay-found}

\paragraph{The explanation we had was impossible.}
The previous version of the suite had one out-of-band number: Saturn's JCT ratio
at $8.75\times$ against a published $6.59\times$. It also had an explanation ---
that first-fit placement fragments Saturn's large multi-node jobs more than the
reference implementation's consolidating placer. That explanation was wrong, and
provably rather than merely unsupported. The replay fleet is single-level: every
node is a direct child of the cluster root, replay jobs carry no locality
constraint and no segment structure, and no crossing penalty is configured, so
\emph{which} free nodes a multi-node job takes cannot change any outcome at all.
The sizes agree: Saturn's largest job is 200 GPUs, and the quota slice carrying
84\,\% of the FIFO-minus-SJF gap tops out at 56 GPUs, with 0.7\,\% of its jobs
larger than one node. It was never a large-job effect.

\paragraph{The real mechanism: sub-node stranding.}
One 1-GPU job landing on an empty 8-GPU node is enough. That node now has seven
free GPUs and an owner, and the allocator treats any node with an owner as
ineligible for a whole-node request --- so a 16-GPU gang cannot touch it. Do
that on enough nodes and the fleet reports plenty of free chips while no gang
can start.

That is what first-fit does here. Saturn is 70.8\,\% single-GPU jobs, 97.1\,\%
in its dominant quota slice. First-fit scans nodes by ascending identifier with
no preference for nodes already partly used, so it opens a fully free node for a
1-GPU job whenever the lower-numbered ones are full, and free capacity
accumulates as 1--7-GPU remainders no job of 8 GPUs or more can claim. The
whole-node rule itself is an approximation whose error was quantified rather
than assumed: Helios's own node count equals $\lceil \mathit{gpus}/8 \rceil$ for
99.66\,\% of windowed multi-GPU jobs, and every exception spreads a job over
\emph{more} nodes than the ceiling, never fewer --- ``never fewer'' being the
half the rule depends on.

Instrumented runs confirmed the mechanism: 100\,\% of measured blocked-idle
chip-seconds come from queue heads needing whole nodes, sub-node heads
contribute 0.0\,\%, and 87.7\,\% of free chips during blocks sit on partly-used
nodes. This happens at an \emph{offered load} of 94.4\,\% --- offered load being
the arriving work as a fraction of the fleet's capacity to do it. Near 100\,\%
queues become extremely sensitive: a few percent of capacity made unusable does
not lengthen the queue by a few percent, it lengthens it by a multiple
\cite{kingman1961}. Those instrumentation figures came from throwaway code not
in the tree --- the diagnostic work that found the mechanism, not numbers the
suite regenerates.

\paragraph{What the fix did and did not buy.}
Switching the sub-node rule to tightest-fit packing --- the \texttt{consolidate}
policy, which fills partly-used nodes before opening empty ones --- closes most
of the gap. The four-cluster mean absolute error against the published ratios
falls from 27.4\,\% to 3.4\,\%, and Saturn's mean FIFO completion time moves
from 75{,}329\,s to 55{,}978\,s against a published 55{,}984\,s. That 0.01\,\%
is luck, not precision --- the tie-break probe below moves the same quantity by
17\,\%. What is supported is the band across four clusters in
Table~\ref{tab:helios-abs}, never a single cell.

\begin{table}[t]
\centering\small
\caption{Absolute quantities: reported, \emph{not} asserted --- continuous
integration does not fail if they move. \texttt{\#Queuing} counts jobs whose wait
exceeds one 60\,s scheduler round. $\ast$ marks the two quantities that moved
\emph{away} from published when the placement model was corrected.}
\label{tab:helios-abs}
\begin{tabular}{lrrrrrr}
\toprule
& \multicolumn{3}{c}{Mean FIFO job completion time (s)} & \multicolumn{3}{c}{\texttt{\#Queuing} jobs} \\
\cmidrule(lr){2-4}\cmidrule(lr){5-7}
Cluster & Published & fleetsim & $\Delta$ & Published & fleetsim & $\Delta$ \\
\midrule
Saturn & 55{,}984 & 55{,}978 & $-0.01$\,\%       & 65{,}991 & 72{,}936 & $+10.5$\,\% \\
Venus  & 64{,}702 & 55{,}714 & $-13.9$\,\%       & 15{,}336 & 13{,}624 & $-11.2$\,\%$^\ast$ \\
Earth  & 19{,}754 & 22{,}463 & $+13.7$\,\%       & 30{,}030 & 32{,}232 & $+7.3$\,\% \\
Uranus & 19{,}758 & 18{,}492 & $-6.4$\,\%$^\ast$ & 16{,}917 & 18{,}449 & $+9.1$\,\% \\
\bottomrule
\end{tabular}
\end{table}

Three qualifications belong with that improvement. One change did not fix
everything: Uranus's absolute FIFO JCT went from $+5$\,\% to $-6.4$\,\%,
overshooting past published rather than approaching it; Venus's queued-job count
drifted from $-9$\,\% to $-11.2$\,\%; and the published \emph{ranking} of
absolute JCTs, which the earlier version happened to reproduce, no longer is ---
though nobody can meaningfully reproduce it, since published Uranus
(19{,}758\,s) and Earth (19{,}754\,s) differ by 0.02\,\%, a dead tie. Only the
ranking of the \emph{ratios} is asserted. Nor are the eight absolute quantities
eight independent confirmations: mean JCT is service plus mean queueing, the
ratio is one mean JCT over another, and the queuing share is queueing over JCT,
so all are functions of the same per-cluster wait distribution and together are
close to one degree of freedom. And the new placer does not dominate the old ---
of the five Saturn quota slices carrying 97\,\% of the gap, three get
\emph{worse} under it, one by 18.5\,\%, the cluster-level win being carried by a
single slice holding 41\,\% of Saturn's jobs. The supported claim is that the policy
reproduces the reference placer's \emph{aggregate} behaviour, never that it is a
better scheduler.

\paragraph{The bands do not discriminate between placement policies.}
\texttt{spread}, a deliberately bad control policy that scatters work onto the
emptiest nodes, was measured across all four clusters and satisfies every
published-facing check above --- both bands on all four, both rank invariants,
the queuing-share ordering --- while being 45\,\% off Saturn's absolute FIFO JCT
and 99\,\% off Earth's. Passing the bands is therefore not by itself evidence
that the placement model is right, and ``in band'' must never be read as
``reproduces the paper''. The four-cluster version of that result is a
documented measurement; what the suite \emph{asserts} is the Saturn instance of
it, so a placement regression fails loudly on one cluster rather than on four.
The point of asserting it at all is that the misreading cannot then be made
silently.

What selects the placer instead, in descending weight, is: the four-cluster mean
absolute ratio error, asserted below 10\,\% and measuring 3.4\,\% against
27.4\,\% for first-fit and 29.5\,\% for the control, since aggregating four
clusters stops any one cluster's ordering noise from deciding it; a mechanism
metric --- the time-averaged count of partly-occupied healthy nodes, fleet state
rather than a queue statistic --- which orders the three policies exactly as
their FIFO completion times do; and Saturn's absolute number, which corroborates
but cannot independently discriminate. It cannot because 35.5\,\% of Saturn's
jobs share an exact submit second, and reversing the tie-break inside those
seconds moves Saturn's FIFO JCT by 17\,\% under first-fit --- enough to invert
which placer looks closer to published. That is a sensitivity probe under a
deliberately wrong arrival order, not an error bar: ascending identifier is
demonstrably the faithful one. Whether the four-cluster discriminator would
survive that perturbation is unmeasured. The probe ran on one cluster under one
scheduler --- Saturn under FIFO, in both tie-break orders --- and extending it
to all four clusters under both schedulers needs 24 more replays. An open
question, not a claim.

\subsection{Closed-form checks}
\label{sec:closed-form}

Trace replay says whether the simulator matches one real cluster; closed-form
checks say whether the engine underneath computes what queueing theory already
knows. The checks form a ladder, from pure algebra at the bottom to a real trace
at the top, and the repository calls each step a \emph{rung}; we keep its word.

Two rungs need their notation unpacked. Queueing theory labels a system by its
arrival process, its service-time distribution and its server count: M/M/c means
memoryless arrivals, memoryless service times and $c$ servers, and M/G/1 means
memoryless arrivals, a general service-time distribution and one server. For a
handful of such systems the mean wait has an exact formula --- Erlang-C for
M/M/c, Pollaczek--Khinchine for M/G/1 --- so we can build the same system in
fleetsim, run it, and check the simulated mean wait against the algebra.

Each rung in Table~\ref{tab:rungs} runs in continuous integration and is chosen
for a specific failure it would catch. Not everything planned is shipped, and
the unshipped items are close cousins of shipped ones, so they are worth naming
precisely. The M/G/1 conservation-law identity --- a statement that a certain
weighted sum of waits is the same under any non-anticipating discipline, and a
different object from the chip-hour accounting invariant in the table --- is a
documented target. So is the M/M/c generalisation of the two-tier priority rung,
whose shipped form is single-server. So are most distributional self-checks on
the traffic generator.

\begin{table}[t]
\centering\small
\caption{Closed-form and invariant rungs, and what each would catch.}
\label{tab:rungs}
\begin{tabular}{p{0.34\textwidth}p{0.56\textwidth}}
\toprule
Rung & What it would catch \\
\midrule
Erlang-C (M/M/c): mean wait and occupancy within 10\,\% of analytic &
A broken queue discipline, a mis-scaled arrival rate, or a double-counting
occupancy integral. \\
Pollaczek--Khinchine (M/G/1) under lognormal service &
Duration-sampler bias the exponential-only rung structurally cannot see, since
the formula depends on the second moment of service time (a distribution's
spread, not just its average, drives the wait). \\
Preemptive-resume priority, two tiers, \emph{one server} (M/M/1), plus a
shielding test under a saturating backlog &
Preemption-accounting errors and leakage of low-priority load into
high-priority wait. Best-effort mean wait is asserted on nothing --- it is
undefined under saturation. \\
Shortest-processing-time optimality, requiring an identical makespan &
A reordering policy that changes how much work is done rather than only when it
is done. \\
EASY backfill \cite{easybackfill1995}, carrying its documented negative case
under honest over-estimates &
An implementation that silently enforces runtime estimates, which real clusters
do not. \\
Invariants: gang atomicity, chip-hour conservation as a per-flush accounting
identity, per-flush allocation
accounting, quota caps with labelled demotions, reservation exclusivity, the
exact-microsecond crossing-penalty identity, byte-identical output for one
scenario and seed on one platform &
The bug class that yields plausible numbers from an internally inconsistent
fleet state. \\
\bottomrule
\end{tabular}
\end{table}

\subsection{What fleetsim deliberately does not reproduce}
\label{sec:anti-goals}

Table~\ref{tab:antigoals} lists the published numbers fleetsim does not try to
match. They are stated so that nobody quotes this paper for something it did not
test; none is asserted anywhere in the suite, so none can surface as a test
failure, which is precisely why they must be written down. Two rows deserve a
note: the fragmentation-delay split and fractional-GPU sharing were both targets
for the release that then spent its budget on placement, and are recorded as
still open rather than quietly re-dated. Separately, the Philly status-split
rung, though written to specification, has never been run against the real
one-gigabyte trace in this build --- the artifact is stored with Git LFS and a
plain fetch returns a pointer file. Published targets, tolerances and structural
assertions on a small synthetic slice exist, but there is no measured Philly
number here, and the suite marks those rows as unmeasured rather than as
agreement.

\begin{table}[t]
\centering\small
\caption{Anti-goals: published results fleetsim deliberately does not
reproduce.}
\label{tab:antigoals}
\begin{tabular}{p{0.31\textwidth}p{0.42\textwidth}l}
\toprule
Published number & Why not & Disposition \\
\midrule
Philly $\sim$52.3\,\% GPU utilization \cite{philly2019} &
A hardware cycle counter, not scheduler occupancy. A discrete-event simulator
produces \emph{allocation} occupancy: a different and higher quantity. &
Out of scope \\
Philly fair-share vs.\ fragmentation-delay split \cite{philly2019} &
Needs per-job delay-cause attribution, which fleetsim does not compute. &
Still deferred \\
Helios QSSF column \cite{helios2021} &
Needs a job-name column absent from the public CSV plus a learned duration
predictor; the SJF-oracle is the reproducible upper-bound proxy. &
Out of scope \\
Alibaba 50\,\% GPU-sharing saving, median 0.042 GPU per instance
\cite{alibabapai2022} &
Needs fractional, sub-chip GPU allocation; fleetsim shares only whole chips
within a node. &
Still deferred \\
Borg absolute occupancy and JCT \cite{borg2015} &
Resources are normalized to $[0,1]$, the data is query-only, and it covers eight
independent 12{,}000-machine cells. &
Never a replay target \\
CPU contention, co-location interference &
Not modelled: an analytical speed model with no memory-bandwidth contention. &
Out of scope \\
\bottomrule
\end{tabular}
\end{table}

\section{Case studies}\label{sec:cases}

Three studies follow, each shipping as a runnable example directory and each a
paired experiment: runs differing by exactly one configuration line. Because
every random draw comes from a stream keyed by name rather than call order,
every other class's arrivals, sizes and durations stay byte-identical across
the pair.

\subsection{Frontier scale: running hot without starving
production}\label{sec:case-frontier}

\texttt{examples/04\_frontier/} is 4 clusters $\times$ 32 pods $\times$ 16
scalable units $\times$ 32 nodes $\times$ 8 chips = 524,288 H100 GPUs, run for
two simulated days with a 120\,s scheduler round at seed 42. Node failures are
off deliberately: a 16,384-node gang is eight times the size of Llama~3's 405B
run, which measured one unplanned interruption every 3.1 hours \cite{llama3}, so
it would break about every 23 minutes (\S\ref{sec:failures} gives the
derivation and the disagreement between the two published anchors), and the
restart churn would drown the scheduling effect. This is therefore a
failure-free result. The two runs differ only in whether a best-effort class
exists: the lowest, always-preemptible priority band, in the sense Borg made
standard \cite{borg2015}. A closed loop holds a fixed number of those jobs in
flight; over the two-day horizon that cycles 11,252 of them through the fleet.

\begin{table}[t]
\centering
\begin{tabular}{lrr}
\toprule
 & No backlog & With backlog \\
\midrule
Occupancy (window) & 0.477 & 0.995 \\
Goodput (window) & 0.966 & 0.880 \\
Productive fraction, occupancy $\times$ goodput & 0.461 & 0.876 \\
\midrule
Eval queue wait, p50 & 61.3\,s & 77.8\,s \\
Batch fine-tune queue wait, p50 & 66.9\,s & 116.4\,s \\
Pretrain queue wait, p50 & 51.7\,s & 163.8\,s \\
Eval queue wait, p99 & 236\,s & 1{,}292\,s \\
Preemptions per minute & 0.068 & 0.565 \\
Jobs finished & 23{,}538 & 34{,}422 \\
\bottomrule
\end{tabular}
\caption{A 524,288-chip fleet with and without a preemptible backlog, seed 42.
The first three rows are measured over the steady-state window (the middle
80\,\% of the run); the rest are full-run figures. The two scopes are not
interchangeable, which is why the example labels them.}
\label{tab:frontier}
\end{table}

Table~\ref{tab:frontier} is the trade. The backlog roughly doubles
steady-state occupancy and finishes 10,884 more jobs. Production medians move
by one or two scheduler rounds --- the largest is the pretrain median, 51.7\,s
to 163.8\,s --- which is the sense in which production waits stay roughly
flat; the tails do not, the eval p99 going from 236\,s to 1,292\,s, roughly 4
minutes to 22. The cost is goodput, 8.6 points of it: the backlog runs with
checkpointing off, which is exactly what makes it instantly reclaimable, so
every eviction discards a whole \emph{stint} --- one contiguous period of
holding an allocation, and with no checkpoint behind it, everything that job
computed since it was last given chips. The product of the two
still nearly doubles. Best-effort jobs are excluded from every wait
distribution by contract, waiting time being meaningless for a closed loop at
saturation.

Then one large job: a 131,072-chip pretraining gang of 32 pod-sized segments,
modelled on the multi-pod shape Llama-3 405B training used \cite{llama3},
submitted at hour 20.57 of both
runs. On the idle fleet it waits 113\,s; above 99\,\% occupancy, 233\,s ---
exactly one 120\,s round more, so the reclaim costs nothing and the wake
boundary costs everything. One wake carries 345 preemptions, allocated chips
dip from 523,946 to 406,406 and recover to 522,998 a round later, and pending
jobs spike from 116 to 441. The gang spans 32 pods in both runs.

\subsection{Occupancy and goodput can move in opposite
directions}\label{sec:case-tradeoff}

\texttt{examples/05\_topology\_tradeoff/} is 8,192 H100s (8 pods $\times$ 16
racks $\times$ 8 nodes $\times$ 8 chips) over seven days at seed 42, failures
off, near 86\,\% utilization so pod-shaped holes are scarce. Pretraining jobs
draw 256, 512 or 1,024 chips and ask for one pod, but softly: they wait ten
minutes for a pod-shaped hole, then take whatever fits. Run~A charges a job
spanning more than one pod a speed multiplier of 0.7, so it needs $1/0.7
\approx 1.43$ times the wall time for the same work; run~B sets that
multiplier to 1.0.

\begin{table}[t]
\centering
\begin{tabular}{lrr}
\toprule
 & A: crossing costs 0.7$\times$ & B: crossings free \\
\midrule
Occupancy (window) & 0.862 & 0.826 \\
Goodput (window) & 0.861 & 0.983 \\
\midrule
Cross-pod placements & 35 of 112 & 17 of 112 \\
Pretrain queue wait, p50 & 484.5\,s & 47.9\,s \\
Pretrain job completion time, p50 & 51{,}002\,s & 43{,}136\,s \\
Pretrains completed & 99 & 103 \\
Eval queue wait, p50 & 31.9\,s & 31.4\,s \\
\bottomrule
\end{tabular}
\caption{Turning the cross-pod speed penalty on raises occupancy and lowers
goodput at the same time (seed 42, 7-day horizon).}
\label{tab:penalty}
\end{table}

Switching the penalty \emph{on} raises occupancy, 0.826 to 0.862, and
\emph{lowers} goodput, 0.983 to 0.861. Both are correct: occupancy is the
fraction of healthy chip-time allocated to some job, goodput the fraction of
that allocated time which produced surviving work. A slowed job holds its
chips about 1.43 times longer for the same work, inflating the first and
deflating the second. The busier-looking run completed four fewer pretrains
and finished them later.

Table~\ref{tab:penalty} also shows the feedback loop that makes this
self-reinforcing, and it reads backwards at first: charging for pod crossings
\emph{doubles} them, from 17 of 112 placements to 35. The chain is that a
penalised job holds its chips 1.43 times longer, so pod-shaped holes are
scarcer when the next pretraining job arrives, so more jobs exhaust their
ten-minute soft wait and take whatever fits. Pretrain queue wait moves
accordingly, 47.9\,s to 484.5\,s.

Two details keep this honest. Run~B's goodput is 0.983 rather than 1.0 because
of checkpoint-save amortization --- a one-hour interval with a 60\,s save
costs about 1.6\,\% --- so the penalty accounts for 12 of those points, not
14. And the effect is confined to large jobs: eval waits are 31.9\,s and
31.4\,s, while the chance of crossing pods rises from 18\,\% at 256 chips to
61\,\% at 1,024. This is the strongest argument available against an
occupancy-only dashboard: occupancy rose 3.6 points and the fleet got worse.

\subsection{Placement: where the small jobs land decides how long the big ones
wait}\label{sec:case-placement}

\texttt{examples/07\_placement\_study/} is deliberately tiny --- 4 racks
$\times$ 8 nodes $\times$ 8 chips = 256 H100s in 32 nodes --- and deliberately
two levels deep, since \texttt{best\_fit} and \texttt{consolidate} are
bit-identical on a single-level fleet. Fourteen days, seed 42, failures off.
The workload is mice (1-, 2- and 4-chip jobs lasting hours) plus gangs of
16, 32 and 64 chips --- 2, 4 and 8 \emph{whole nodes} --- at an offered load of
78\,\% by construction, 81\,\% realised at this seed, because Poisson arrivals
and a lognormal duration tail land above the mean here. A best-effort FIFO
scheduler holds ordering fixed, so only the
placement policy varies. It reproduces the mechanism the Helios replay
identified: a 1-chip job landing on a completely free node converts eight
whole-node chips into a remainder no gang can claim, and the gangs, not the
mice, pay for it.

\begin{table}[t]
\centering
\begin{tabular}{lrrrr}
\toprule
 & \texttt{first\_fit} & \texttt{best\_fit} & \texttt{consolidate} & \texttt{spread} \\
\midrule
Stranded whole nodes, mean (of 32) & 6.55 & 5.31 & 5.37 & 25.31 \\
64-chip gang queue wait, mean (s) & 19{,}056.8 & 16{,}195.3 & 16{,}314.7 & 739{,}842.1 \\
64-chip gangs ever started (of 109) & 107 & 107 & 107 & 29 \\
Mice queue wait, p50 (s) & 31.0 & 31.4 & 31.3 & 30.1 \\
Occupancy (full run) & 0.8141 & 0.8141 & 0.8141 & 0.5829 \\
Jobs completed & 3{,}436 & 3{,}436 & 3{,}436 & 3{,}309 \\
\bottomrule
\end{tabular}
\caption{Four placement policies on identical hardware, identical arrivals and
identical ordering (seed 42, 14-day horizon). A ``stranded'' node is a healthy
node that is partly occupied, so its free chips cannot satisfy any whole-node
request. The whole study runs in about 10\,s on a laptop: about 1\,s each for
the three packing arms and about 7\,s for \texttt{spread}, whose queue is
roughly 30$\times$ deeper.}
\label{tab:placement}
\end{table}

Tightest-fit packing (\texttt{best\_fit}; \texttt{consolidate} is within noise
of it here) cuts time-average stranded nodes from 6.55 to 5.31 of 32
and the mean queue wait of an 8-node gang from 19,056.8\,s to 16,195.3\,s
(Table~\ref{tab:placement}). The mice pay for it: their median wait is flat,
31.0 against 31.4\,s, but their mean rises from 80.7\,s to 95.3\,s --- about
15 seconds of mouse wait buying about 2,900 seconds of gang wait.
\texttt{spread} prices the mechanism from the other side: it strands 25.31 of
32 nodes, never starts 80 of the 109 8-node gangs and delivers 28.4\,\% fewer
chip-hours on identical hardware --- and it averages only the gangs it did
start.

The negative result matters more. Full-run occupancy, goodput and allocated
chip-hours are \emph{bit-identical} across \texttt{first\_fit},
\texttt{best\_fit} and \texttt{consolidate} --- 0.8140502977503281,
0.9831464030337717 and 70,021.35041129222 chip-hours --- because the same
3,436 jobs complete for the same durations in all three, and windowed
occupancy moves 0.2 percentage points. The entire placement effect is
\emph{when} work ran, never how much: a placement study reporting only
occupancy has measured almost nothing.

Three caveats travel with it. At seed 1234 the mechanism metric still improves
by about 20\,\% while the outcome regresses: mean gang wait 8,205\,s to
9,407\,s, 15\,\% worse. A hard \texttt{within: rack} constraint inverts the
ranking outright --- 64-chip mean wait 39,677\,s under \texttt{first\_fit}
against 121,486\,s under \texttt{best\_fit} --- because node-level
tightest-fit spreads gangs across racks, so no rack ever drains. And this
scheduler is best-effort, so a stuck gang does not hold up the mice behind it:
a floor, not a best case. ``\texttt{best\_fit} is a better scheduler'' is not
a claim this study supports.

\section{Implementation, testing and reproducibility}\label{sec:impl}

\paragraph{Language and size.} fleetsim is Python 3.11 or newer under
Apache-2.0, depending at runtime on numpy, pandas, pyarrow, pyyaml and
optionally matplotlib. The event loop is hand-written rather than built on a
simulation framework --- full control over tie-breaking and determinism is
worth about 200 lines --- and the local web application is pure
standard-library \texttt{http.server}, so running it adds no dependency. The
simulator is 22,563 lines of Python across 47 files; the tests are 18,896
lines across 36 files plus 2,585 lines of validation rungs across 12, and
pytest collects 1,024 of them. The four placement policies behind
Section~\ref{sec:case-placement} are 340 lines together; the FIFO scheduler is
48.

\paragraph{Determinism, and its per-platform limit.} Simulated time is a
64-bit integer count of microseconds, so ordering and tie-breaking are exact.
Random draws come from numpy seed sequences whose spawn key is the first 16
bytes of the SHA-256 of the stream's \emph{name}, never Python's salted
\texttt{hash()}, so a
given stream yields the same numbers in every process on every machine. The
contract is that an identical scenario and seed produce byte-identical Parquet
and JSON, on a given platform, and a CI test enforces it by running the whole
pipeline twice with failures, maintenance and preemption switched on.

That qualifier is load-bearing. The lognormal and exponential samplers call
the platform's C math library, whose \texttt{exp} and \texttt{log} differ in
the last bit between glibc and Apple's implementation. An early release (v0.1;
the project is at v0.8) shipped a compatibility test pinning a SHA-256 of the
arrival stream captured on macOS, and it failed the first time CI ran on Ubuntu. The replacement compares golden
\emph{data} --- the first 1,000 arrivals at nine significant digits, plus
aggregate counts --- with tolerances: 2 microseconds on times, $10^{-6}$
relative on durations, and at most 5 differing rows per 1,000 for draws
sitting on a discrete knife edge. Cross-platform reproduction means the same
distributions, not the same bytes.

\paragraph{Continuous integration, and what a clean machine caught.} One
GitHub Actions workflow runs on every push and pull request: Ubuntu, a Python
3.11 and 3.12 matrix, an editable install of the package and the out-of-tree
plugin example, \texttt{pytest tests validation -q}, then a smoke test running
\texttt{fleetsim validate} and \texttt{fleetsim run} on the minimal example.
It runs with no network in about 80\,s with exactly four skips --- two Helios
full-trace rungs, one Philly, one frontier-budget gate --- and anything else
skipping is a signal.

Two defects surfaced only because CI is a clean room. The first was an
environment leak: a test asserting that every shipped example validates passed
locally only because the developer's virtual environment still had the
example-03 plugin scheduler installed from an earlier session, and that
example names an out-of-tree scheduler which enters the registry only when its
package is installed. The fix skips an example whose scheduler is unregistered
--- validation \emph{should} reject unknown names --- while asserting a floor
on how many examples were actually checked, so the skip cannot hollow the test
out; CI now installs the plugin. The second was a concurrency deadlock a
laptop could not reach: the 3.12 job hung at 91\,\% and was killed by the job
timeout while 3.11 passed the same commit, the developer machine having many
cores and the hosted runners two. The cause was a helper process that starts
lazily, combined with a setup step that re-ran on every worker-pool creation and
mutated interpreter-global state; on a two-core runner the two raced. Building
that context once per process fixed the hang, and a faulthandler timeout was
added so that the next hang of any kind would fail loudly with every thread's
stack. That dump then named a deeper bug: CPython invokes a future's
done-callback while holding the executor's shutdown lock, so a callback
calling \texttt{shutdown()} deadlocks itself. A broken pool is now disposed on
a separate daemon thread.

\paragraph{Adversarial review.} Several releases were reviewed adversarially ---
a pass whose goal was to find defects rather than approve --- and each such pass
leaves a regression file that lists its findings:
\texttt{test\_\allowbreak review\_\allowbreak fixes.py} in \texttt{tests/}, plus
its \texttt{v02}, \texttt{v04}, \texttt{v05} and \texttt{v05\_hardening}
successors. Five files across eight releases, so this is a habit rather than a
gate. The same
discipline governs the published figures: each lives once, as data, in a
single module from which the validation tests import their assertion constants
and which the web application serves verbatim, and a test checks by object
identity that nobody has retyped a number into a test.

\paragraph{Reproducing the validation.} The offline checks are one command,
\texttt{pytest tests validation -q}, and \texttt{fleetsim validation run}
prints the vendored-slice checks in about 5\,s. The trace results are opt-in
and environment-gated because they download real data:
\texttt{FLEETSIM\_HELIOS\_FULL=1} \texttt{pytest -m trace\_full}
\texttt{validation/test\_helios\_ratio.py} fetches the 36\,MB Helios archive
\cite{helios2021} and replays four clusters under two schedulers in 266\,s.
The Philly rung \cite{philly2019} is gated the same way but needs \texttt{git
lfs pull} first, since a plain fetch of that 1\,GB artifact yields a 135-byte
pointer; it was written to specification and has not been run against the real
trace in this build. Each rung self-skips unless its variable is set and never
downloads implicitly. \texttt{fleetsim validation cite} prints each trace's
required attribution.

\paragraph{Tooling.} Six subcommands cover the loop: \texttt{fleetsim
validate}, \texttt{run}, \texttt{compare}, \texttt{plot}, \texttt{viz} and
\texttt{serve}. Two properties of them carry an argument rather than a feature
list. \texttt{viz} renders a run directory into one self-contained HTML replay
that makes zero external requests, and a test greps both the template and the
output for any external URL or fetch call, so the claim is enforced rather than
intended. And parameter sweeps run one simulation per grid cell in worker
processes: the same sweep run twice is byte-identical cell for cell, so
parallelism does not weaken determinism, and a test enforces that too.

\section{Related work}
\label{sec:related}

Recall from \S\ref{sec:intro} that simulators are best compared by
\emph{altitude}: what one simulated run
represents. A run of \texttt{fleetsim} represents weeks of a whole fleet ---
many clusters, many tenants, thousands of jobs arriving, queueing, running,
being preempted and occasionally dying with their nodes. Most existing tools sit
at a different altitude, which is not a deficiency in them: a tool that models
one job's network traffic in microseconds is not a worse fleet simulator, it is
not a fleet simulator at all, and it answers questions \texttt{fleetsim} cannot.
Table~\ref{tab:related} summarises the landscape. One caveat: this is a survey of
what each tool's papers and documentation describe, not a benchmark --- we did
not install, run or measure any of these systems.

\begin{table}[!t]
\centering
\small
\caption{Related tools by altitude. Characterisations are taken from each
tool's published description, not from measurements we made.}
\label{tab:related}
\begin{tabular}{@{}p{2.6cm}p{2.2cm}p{4.0cm}p{4.3cm}@{}}
\toprule
\raggedright Family & \raggedright Altitude & \raggedright What it does well & \raggedright What a fleet question needs that it omits \tabularnewline
\midrule
\raggedright Blox \cite{blox2024}; the simulators inside Gavel, Pollux, Sia
  & \raggedright One cluster
  & \raggedright Fast, faithful replay of published deep-learning scheduling policies; Blox reports seven schedulers in 12--1157 lines each
  & \raggedright Capacity is a flat GPU count: no topology levels, no node failures or repair, no inference, no multi-cluster \tabularnewline
\addlinespace
\raggedright ASTRA-sim \cite{astrasim2023}, SimAI \cite{simai2025}
  & \raggedright One job $\times$ network
  & \raggedright Collective communication and per-iteration time at packet and flow detail
  & \raggedright There is no queue: the job is already placed, so there are no arrivals, no waiting and no preemption \tabularnewline
\addlinespace
\raggedright Vidur \cite{vidur2024}, LLMServingSim
  & \raggedright One inference service
  & \raggedright Request batching, KV-cache behaviour and tail latency for LLM serving
  & \raggedright Inference only, one service; training gangs, quota and failures are outside the model \tabularnewline
\addlinespace
\raggedright Batsim \cite{batsim2016}, the UBCCR Slurm simulator
  & \raggedright HPC batch
  & \raggedright Decades of batch-queue semantics, backfill \cite{easybackfill1995}, and a clean engine/policy separation
  & \raggedright No accelerator cost model: preemption is not priced in lost training work, and there is no checkpoint or gang-restart notion \tabularnewline
\addlinespace
\raggedright AIReSim \cite{aireesim}
  & \raggedright Reliability of one job
  & \raggedright Failure, repair and restart accounting for a single long training run
  & \raggedright No scheduler, no queue, no competing tenants \tabularnewline
\addlinespace
\raggedright kube-scheduler-simulator \cite{kubescheduler}
  & \raggedright Kubernetes placement
  & \raggedright Explains, decision by decision, why a pod landed where it did
  & \raggedright Wall-clock, no simulated time: no job completion time, no occupancy over a horizon \tabularnewline
\addlinespace
\raggedright Borg \cite{borg2015,borgtng2020}, MAST \cite{mast2024}, Singularity \cite{singularity2022}
  & \raggedright Fleet --- the existence proof
  & \raggedright Run real fleets at exactly this altitude and publish what the hard problems are
  & \raggedright They are production systems, not simulators; there is no public artefact in which to run an experiment \tabularnewline
\bottomrule
\end{tabular}
\end{table}

\paragraph{Single-cluster scheduler simulators.}
Blox \cite{blox2024} and the simulators built inside scheduling papers such as
Gavel, Pollux and Sia take a job trace and a policy, apply them to one cluster
of $N$ interchangeable GPUs, and report completion times. They are the closest
prior art, and \texttt{fleetsim} borrows from them directly: Blox's decomposition
of a scheduler into an \emph{ordering} policy (who runs next) and a
\emph{placement} policy (where it goes) is the API we adopted, and its convention
of measuring the steady middle of a run is our metrics window. What a flat GPU
count cannot express is a job failing to fit for structural reasons. Our
placement study makes the consequence concrete: on a 256-chip fleet at one seed,
changing only the placement rule left occupancy identical to the last digit
(0.8141 under three of four policies) while the mean queue wait of eight-node
gangs moved from 19{,}057\,s to 16{,}195\,s. A simulator whose capacity is an
integer cannot represent that mechanism; node failures are the second gap.

\paragraph{Per-job network simulators.}
ASTRA-sim \cite{astrasim2023} and SimAI \cite{simai2025} simulate one training
job's iteration in detail --- collectives, fabric topology, congestion. This is
precisely the layer \texttt{fleetsim} collapses into a single number: a
placement spanning a configured topology level multiplies the job's progress
rate by a configured multiplier, and there is no default multiplier because the
honest value is fleet-specific. The two altitudes compose rather than compete ---
a network-level study is how a lab would derive that multiplier. What these tools
do not have is a queue: one already-placed job has nothing to wait behind.

\paragraph{Inference-serving simulators.}
Vidur \cite{vidur2024} and LLMServingSim model a serving deployment: request
arrivals, batching, KV cache, latency percentiles. \texttt{fleetsim} models
inference at replica granularity --- a replica being a fixed block of chips held
for a service --- and stops there, because per-request burstiness sits below
fleet altitude.

\paragraph{HPC batch simulators.}
Batsim \cite{batsim2016} and the Slurm simulator carry decades of batch-queue
semantics, and \texttt{fleetsim} copies Batsim's central architectural choice:
the engine mutates state, policies emit intents, and the engine validates and
applies them, in-process rather than over a socket. EASY backfill
\cite{easybackfill1995} is a shipped \texttt{fleetsim} scheduler for the same
reason. The gap is not gangs as such --- MPI jobs are gangs --- but the
accelerator cost model around them: killing a training job destroys the work done
since its last checkpoint, a GPU job cannot be suspended in place because its
memory is held, and a restart costs minutes. Those costs are what make preemption
policy interesting, and what these models do not price.

\paragraph{Production fleet managers.}
Borg \cite{borg2015,borgtng2020}, MAST \cite{mast2024} and Singularity
\cite{singularity2022} operate at fleet altitude, and supply much of
\texttt{fleetsim}'s substance: Borg's priority bands are the tier semantics, and
published reliability figures \cite{metareliability2024,llama3} set the failure
defaults. Borg TNG's tail of per-job \emph{resource-hours} ($\alpha \approx
0.69$--$0.77$, the top 1\,\% of jobs accounting for 97--99\,\% of load) is a
target the traffic generator is meant to satisfy --- not a configured parameter,
and not yet checked by a shipped rung. What
none of them ships is something you can run, and the public traces
\cite{philly2019,helios2021,alibabapai2022,acme2024} record only what happened
under one policy, never what a different one would have done. That counterfactual
is the gap. One precision about the borrowing: \texttt{fleetsim}'s quota model is
inspired by MAST but is not a model \emph{of} it (Section~\ref{sec:limitations}).

The workload generator likewise stands in the workload-modelling tradition of
rigid parallel jobs \cite{feitelson2003}, with non-Poisson arrivals
\cite{paxsonfloyd1995} and lognormal duration bodies taken from the Helios and
Acme characterizations \cite{helios2021,acme2024} --- and the closed-form checks
of
Section~\ref{sec:validation} are ordinary queueing theory \cite{kingman1961}. A
cross-simulator check against Blox's published percentiles was planned and never
built.

\section{Limitations and future work}
\label{sec:limitations}

\paragraph{No fractional GPUs.}
The smallest unit \texttt{fleetsim} allocates is one whole chip. Alibaba PAI's
headline results \cite{alibabapai2022} --- roughly a 50\,\% saving from GPU
sharing, with a median instance requesting 0.042 of a GPU --- are therefore out
of reach, and no placement policy substitutes for the missing mechanism. It was
a target of the previous release that the release did not reach, and we record
it as still open rather than quietly re-dating it.

\paragraph{No per-job delay-cause attribution.}
\texttt{fleetsim} records how long each job waited, not \emph{why}. Philly's
split of queueing delay into a fair-share component and a fragmentation component
\cite{philly2019} --- the result that motivated this design, where fragmentation
accounts for 59--98\,\% of multi-GPU delay --- cannot be reproduced, and is a
second promise carried over from the release before last, still open. The metric
we did add, a count of
partially-occupied nodes, is a down payment, not a substitute, and it is
directional rather than proportional: an 8.9\,\% reduction in its
time-average accompanied a 25.7\,\% reduction in mean FIFO completion time,
because a mean over all instants dilutes exactly the moments when the queue head
needs a whole free node.

\paragraph{The performance model is analytical.}
A job's progress rate is the product of configured per-level multipliers. The
identity \emph{completion $=$ start $+$ overhead $+$ remaining work / rate} is
pinned exactly in integer microseconds, but the rate is an input, not a
derivation, and there is no default crossing penalty. Chip heterogeneity is worse
off: the throughput hook always returns 1.0, and scenarios using unpinned chip
types are rejected outright. CPU contention and co-location interference are not
modelled at all.

\paragraph{Failures are sampled from a rate, not from a physical model.}
One exponential draw at the fleet-wide aggregate rate, a thinning test, then a
uniform victim pick --- at a default mean time between failures of 42 node-days
with a fixed cause mix (60\,\% GPU/HBM, 10\,\% network, 13\,\% software, 17\,\%
other). Nothing in that is physical: no correlated failures (a rack losing power
takes one node at a time here), no thermal or age dependence, and unreliable
``lemon'' nodes are off by default. Spot preemption is deliberately harsher than
the products it abstracts --- zero notice, no banked progress --- so spot losses
are an upper bound, not an estimate.

\paragraph{One trace family, validated in depth.}
The quantitative validation is Helios \cite{helios2021}: four clusters, both
schedulers, the full September window. Philly \cite{philly2019} is exercised only
at converter level, and even there the full check, though written, was never run
against the real one-gigabyte artefact, so every Philly row reports
``unmeasured'' rather than a number; only structural assertions on a synthetic
2{,}000-row slice run in CI. Alibaba PAI is registered but never replayed. One
trace family is one lab's workload, one scheduler lineage, one hardware
generation.

\paragraph{The absolute numbers are order-sensitive.}
On Saturn, 35.5\,\% of the 105{,}969 jobs in the window share an exact submit
second, and 92.9\,\% share a 60-second scheduler round --- so the order in which
tied arrivals are served is a modelling choice, not a property of the trace.
Flipping the FIFO tie-break from ascending to descending job id moves Saturn's
mean FIFO completion time from 75{,}329\,s to 62{,}549\,s, a 17.0\,\% move, and
inverts which placement policy looks closer to published. Ascending id is
demonstrably the faithful order (ids are assigned at submission and are monotone
in submit time on three of the four clusters), so this is a sensitivity probe
under a deliberately wrong order, not an error bar. Its magnitude still matters:
17\,\% is wider than the $\pm 14$\,\% agreement we report on absolute completion
times, and a large fraction of the 25.7\,\% move the placement change itself
produced. No single cluster's absolute number can select a modelling choice,
which is why our discriminator aggregates four clusters --- and whether that
discriminator survives the same perturbation is \emph{unmeasured}: the probe ran
on one cluster under one scheduler (Saturn under FIFO, in both tie-break
orders), and extending it to all four clusters under both schedulers needs 24
more replays.

\paragraph{Single-metro scheduling only.}
There is one scheduling stage over one metro. Two-stage scheduling across metros
--- the MAST shape \cite{mast2024}, where a first stage decides which region a job
belongs in --- is schema-present and rejected as not implemented, as are TPU
optical-switch predicates, multi-gang jobs and autoscaling inference. The quota
model is also simpler than the systems that inspired it: \texttt{fleetsim}
commits quota at admission against queued demand and demotes irreversibly,
whereas MAST and comparable systems evaluate against running usage at scheduling
time and treat over-quota work as borrowed capacity that can be handed back. The
quota scenario shipped as \texttt{examples/06\_economics/} is not among the case
studies of \S\ref{sec:cases}; its demotion counts are partly an artefact of that
admission-time choice, which is why we do not quote them here.

\paragraph{Smaller ones, stated plainly.}
The backfill scheduler implements the conservative one-rule form of EASY, so
measured backfill rates undershoot published EASY reproductions. Byte-identical
output holds only within a platform, since the floating-point samplers differ in
the last bit across operating systems. And the performance envelope in the design
document --- six simulated months of a 100{,}000-GPU fleet --- is a target, not a
measurement: what has been measured is six months of a 2{,}048-chip fleet in
about nine minutes, and two days of a 524{,}288-chip fleet in about eighty
seconds, on a laptop.

\paragraph{Future work.}
Five items, in priority order. \emph{Fractional GPUs}, the single
change that unlocks a whole trace family's published results. \emph{Per-job
delay-cause attribution}, which would reproduce Philly's split and turn the
stranding metric into a per-job account of where a wait came from. \emph{More
trace families} --- the Philly artefact, then Alibaba PAI --- because one
validated family is the weakest part of Section~\ref{sec:validation}. \emph{TPU
torus geometry}: TPU pods wire their chips into a three-dimensional torus --- a
grid whose edges wrap around --- and a job is given a rectangular sub-block of
it, a \emph{slice}. Feasibility then depends on which free \emph{cubes}
(contiguous $4 \times 4 \times 4$ blocks) exist rather than on how many chips are
free, and the published optical-switch result removes contiguity constraints for
cube-multiple slices \cite{tpuv4isca2023}. And a
\emph{throughput matrix for heterogeneous chips}, replacing the hook that returns
1.0, so a fleet mixing chip generations can be modelled rather than approximated
as uniform.

\section{Conclusion}
\label{sec:conclusion}

We built an open-source discrete-event simulator that operates at fleet altitude:
a hierarchy of clusters, pods, racks and nodes; gang scheduling with atomic,
immutable allocations; priority preemption priced in lost work; node failures,
repair and maintenance drains; and a generative traffic model drawn from
published trace characterisations.

We then tried to break it. Against the four Helios clusters \cite{helios2021},
\texttt{fleetsim} reproduces the published ratio between FIFO and
shortest-job-first mean completion times with a mean absolute error of 3.4\,\%,
and lands within $\pm 14$\,\% of the published absolute completion times --- a
weaker, reported-not-asserted result the repository deliberately calls ``right
ballpark'' rather than reproduction. Both hold only under three modelling
choices we name every time we state the result, and inside tolerance bands wide
enough that a deliberately bad placement policy passes them too. The engine also
reproduces four textbook queueing results it could have failed --- Erlang-C,
Pollaczek--Khinchine, two-class preemptive-resume priority and
shortest-processing-time optimality --- plus a behavioural property test of EASY
backfill. And we have written down what it does not reproduce, and why.

If one sentence survives: matching a published table was not the useful part ---
the useful part was that matching it forced three modelling choices into the
open, and that occupancy, the metric everyone reports, was bit-for-bit identical
across placement policies whose queueing behaviour was not.


\bibliographystyle{plain}
\bibliography{refs}

\end{document}
